\section{Introduction}
\label{sec:intro}
% The authors' introduction (Nov 20), with the corrections accepted.
% Keep editing it in your own words.

Textured products such as carpet, leather, wood and tile are judged by their
appearance, so even a small defect can make them unsellable. Industry needs
computers and software to do visual inspection and to check every product on
the production line. Companies mostly produce good products, so if we wait to
collect hundreds of images of defective products, it takes a long time. Also,
inspection by human workers is slow, and they miss small defects when they
become tired.

Defects are also varied, and labelling is very expensive, because the exact
defect must be drawn pixel by pixel. Instead of learning what defects look
like, methods therefore learn what normal products look like and flag anything
different. This approach is called anomaly detection (AD), and the MVTec AD
dataset is its standard benchmark~\cite{bergmann2021mvtec}.

Training-free AD methods differ mainly in where they find their reference for
``normal''. Zero-shot single-image methods use no reference images at all: they
only compare each patch of the test image with the rest of the same image and
flag regions that break its homogeneity~\cite{aota2023zeroshot,ardelean2024highfidelity,yao2025grnr,zhang2025nlpa}.
Few-shot methods build a memory bank from the patches of a few defect-free
reference images and give each test patch a score from its distance to the
nearest entry~\cite{roth2022patchcore,damm2025anomalydino,lendering2026subspacead}.

Batched zero-shot methods compare patches across many unlabelled test
images~\cite{li2024musc,legia2025consistent}. Frozen self-supervised features,
in particular DINOv2~\cite{oquab2024dinov2}, make all three families practical
without any training.

Single-image methods depend on an assumption: most of the image is normal.
Aota et~al.~\cite{aota2023zeroshot} state this assumption for pixel-level
detection, and the authors of feature correspondence analysis
(FCA)~\cite{ardelean2024highfidelity} acknowledge a limitation when anomalies
extend across most of the image. The underlying assumption is therefore
expected to become less reliable as defect size increases, because a large
defect becomes part of the reference that each patch is compared with.
Published results, however, are reported as averages over all defects, which
hides any dependence on size. Defect size has been studied for a single method
and small defects~\cite{lehr2024limitations}, and contamination of the
reference set of few-shot detectors has been studied~\cite{villanueva2026cleaning},
but to our knowledge no study has measured how single-image self-similarity and
a few-shot memory bank compare across defect sizes on the same frozen features,
or tested whether the defect contaminating the image's own reference explains
the difference.
% NOTE: confirm "to our knowledge" with Google Scholar and arXiv
% (strings G1--G3 in the workbook) before submission.

This study fills that gap with four research questions:
\begin{enumerate}
  \item[RQ1:] How much does a four-image memory bank improve on a simple
  mean-similarity single-image baseline?
  \item[RQ2:] How does that improvement depend on defect size?
  \item[RQ3:] How much of the size dependence comes from the defect
  contaminating the image's own reference?
  \item[RQ4:] Does a robust published single-image method (FCA) show the same
  dependence?
\end{enumerate}

We do not propose a new detector. We compare a mean-similarity single-image
baseline (Route~A), our implementation of the memory-bank idea of
AnomalyDINO~\cite{damm2025anomalydino} (Route~B) and the official FCA code on
the five MVTec AD textures, with the same frozen DINOv2 features and the same
metric code.

The main contributions of this work are as follows:
\begin{enumerate}
  \item A controlled comparison of references. On the same frozen DINOv2
  features and with the same metric code, we show that a four-image memory bank
  outperforms mean self-similarity on all five textures (pixel-level average
  precision, Pixel AP, 32.1\% vs.\ 50.0\%), and that a robust published
  single-image method (FCA) outperforms both on segmentation.
  \item A size analysis. On 382 defective images, we show that the advantage of
  the memory bank grows with defect size (a slope difference of $-2.72$ points
  of per-image area under the ROC curve, AUROC, per unit of log defect area),
  while the slope of FCA cannot be distinguished from that of the memory bank;
the strong size dependence is therefore not shared by every single-image
method.
  \item An intervention study. On real and synthetic defects, we show that
  removing the defect from the image's own reference recovers about half of the
  size slope on real defects and about two thirds on synthetic defects, whereas median and trimmed-mean references recover little; on real defects
the rest of the slope remains unexplained.
\end{enumerate}
We also report the ideas that did not work (routing between the two methods, an
image-level mean score and feature drift), and we fixed every success criterion
before running the analysis; code and results are kept in a tagged repository that will be made public
on publication.

The rest of this paper is organised as follows.
Section~\ref{sec:related} reviews the related literature and
Section~\ref{sec:methods} describes the methodology; Section~IV presents the
results, Section~V discusses them and Section~VI concludes.
% When Results, Discussion and Conclusion are switched on in main.tex,
% replace IV, V, VI with \ref{sec:results}, \ref{sec:discussion}, \ref{sec:conclusion}.
